\documentclass[12pt, a4paper]{article}

% ── Packages ──────────────────────────────────────────────────────
\usepackage{fontspec}
\setmainfont{Sylfaen}
\usepackage{amsmath, amssymb, amsthm}
\usepackage{mathtools}
\usepackage{enumitem}
\usepackage{tikz}
\usetikzlibrary{calc, positioning, arrows.meta, fit, patterns, decorations.pathreplacing}
\usepackage{geometry}
\geometry{margin=2.4cm}
\usepackage{xcolor}
\usepackage{tcolorbox}
\tcbuselibrary{theorems, skins, breakable}
\usepackage{booktabs}
\usepackage{hyperref}
\hypersetup{colorlinks=true, linkcolor=blue!60!black, urlcolor=blue!60!black}

\setlength{\parindent}{0pt}
\setlength{\parskip}{6pt}

% ── Colours ───────────────────────────────────────────────────────
\definecolor{setA}{RGB}{66, 133, 244}   % blue
\definecolor{setB}{RGB}{234, 67, 53}    % red
\definecolor{setC}{RGB}{251, 188, 4}    % yellow
\definecolor{theoryBg}{RGB}{232, 245, 253}
\definecolor{theoryFrame}{RGB}{66, 133, 244}
\definecolor{stmtBg}{RGB}{255, 243, 224}
\definecolor{stmtFrame}{RGB}{230, 126, 34}
\definecolor{ansBg}{RGB}{232, 250, 232}
\definecolor{ansFrame}{RGB}{39, 174, 96}

% ── Environments ──────────────────────────────────────────────────
\newtcolorbox{theorybox}[1][]{%
  colback=theoryBg, colframe=theoryFrame,
  fonttitle=\bfseries, title={\small Տեսության ամփոփում. #1},
  breakable, enhanced, arc=2mm,
  left=5pt, right=5pt, top=5pt, bottom=5pt
}

\newtcolorbox{problembox}{%
  colback=stmtBg, colframe=stmtFrame,
  fonttitle=\bfseries, title={\small Խնդրի դրվածք},
  breakable, enhanced, arc=2mm,
  left=5pt, right=5pt, top=5pt, bottom=5pt
}

\newtcolorbox{answerbox}{%
  colback=ansBg, colframe=ansFrame,
  arc=2mm, left=5pt, right=5pt, top=3pt, bottom=3pt
}

\newcommand{\Z}{\mathbb{Z}}
\newcommand{\R}{\mathbb{R}}
\newcommand{\N}{\mathbb{N}}
\newcommand{\powset}[1]{2^{#1}}
\newcommand{\comp}[1]{\overline{#1}}
\newcommand{\symdiff}{\mathbin{\triangle}}

% ── Title ─────────────────────────────────────────────────────────
\title{\textbf{Գործնական 5 — Կոմբինատորիկա (հիմնական կանոններ)}\\[4pt]
       \large Ընտրված խնդիրների լուծումներ\\[2pt]
       \normalsize \S1: Խնդիրներ 3, 5, 10 \quad \S2: Խնդիրներ 5, 10, 12}
\author{Դիսկրետ մաթեմատիկա — ԵՊՀ}
\date{}

\begin{document}
\maketitle
\tableofcontents
\newpage

% ══════════════════════════════════════════════════════════════════
\part*{Բաժին 1. Հիմնական կանոններ (Գումարի և Արտադրյալի կանոններ)}
\addcontentsline{toc}{part}{Բաժին 1. Հիմնական կանոններ}
% ══════════════════════════════════════════════════════════════════

% ══════════════════════════════════════════════════════════════════
\section{Խնդիր 3 — Տարբեր առարկաների երկու գրքի ընտրություն}
% ══════════════════════════════════════════════════════════════════

\begin{theorybox}[Գումարի և Արտադրյալի կանոններ]
Երկու ամենահիմնարար հաշվարկման սկզբունքները․

\begin{itemize}[nosep]
  \item \textbf{Գումարի կանոն (Այլընտրանքների կանոն).}
    Եթե գործողությունը կարող է կատարվել $k$ փոխադարձաբար բացառող եղանակներից մեկով,
    և $i$-րդ եղանակը կարող է իրականացվել $n_i$ ձևով, ապա ձևերի ընդհանուր քանակը
    $n_1 + n_2 + \cdots + n_k$ է։

  \item \textbf{Արտադրյալի կանոն (Հաջորդական ընտրությունների կանոն).}
    Եթե գործողությունը բաղկացած է $k$ հաջորդական քայլերից, որտեղ $i$-րդ քայլը կարող է
    կատարվել $n_i$ ձևով (անկախ մյուս քայլերի ընտրություններից),
    ապա ձևերի ընդհանուր քանակը $n_1 \cdot n_2 \cdots n_k$ է։
\end{itemize}

\medskip
\textbf{Համակցված օգտագործում.}\;
Երբ խնդիրը կարող է տրոհվել մի քանի \emph{դեպքերի} (գումարի կանոն),
որոնցից յուրաքանչյուրը ներառում է \emph{հաջորդական ընտրություններ} (արտադրյալի կանոն),
մենք հաշվում ենք յուրաքանչյուր դեպքի քանակը արտադրյալի կանոնով, ապա գումարում ենք արդյունքները։
\end{theorybox}

\begin{problembox}
Քանի՞ եղանակով կարելի է ընտրել տարբեր առարկաների վերաբերյալ երկու գիրք,
եթե կան 15 գիրք Ինֆորմատիկայից, 12 գիրք Մաթեմատիկայից և 10 գիրք Ֆիզիկայից։
\end{problembox}

\subsection*{Լուծում}

Քանի որ երկու ընտրված գրքերը պետք է լինեն \emph{տարբեր} առարկաներից,
մենք դիտարկում ենք առարկաների բոլոր հնարավոր զույգերը։
Կա երեք դեպք (գումարի կանոն), և յուրաքանչյուր դեպքում կիրառում ենք
արտադրյալի կանոնը․

\begin{center}
\renewcommand{\arraystretch}{1.4}
\begin{tabular}{lccc}
\toprule
\textbf{Դեպք} & \textbf{Առարկա 1} & \textbf{Առարկա 2} & \textbf{Քանակ} \\
\midrule
Ինֆորմատիկա + Մաթեմատիկա & 15 ընտրություն & 12 ընտրություն & $15 \cdot 12 = 180$ \\
Ինֆորմատիկա + Ֆիզիկա     & 15 ընտրություն & 10 ընտրություն & $15 \cdot 10 = 150$ \\
Մաթեմատիկա + Ֆիզիկա     & 12 ընտրություն & 10 ընտրություն & $12 \cdot 10 = 120$ \\
\bottomrule
\end{tabular}
\end{center}

Ըստ գումարի կանոնի (երեք դեպքերը փոխադարձաբար բացառող են), ընդհանուր քանակն է․
\[
  15 \cdot 12 \;+\; 15 \cdot 10 \;+\; 12 \cdot 10
  \;=\; 180 + 150 + 120
  \;=\; 450.
\]

\begin{answerbox}
\centering
Տարբեր առարկաների երկու գիրք ընտրելու $\boxed{450}$ եղանակ կա։
\end{answerbox}

\textbf{Այլընտրանք (լրացման մեթոդ).}\;
Ընդհանուր զույգերի քանակը բոլոր $15+12+10=37$ գրքերից՝ $\binom{37}{2}$։
Հանենք նույն առարկայի զույգերը՝ $\binom{15}{2} + \binom{12}{2} + \binom{10}{2}$։
Արդյունքը՝ $\binom{37}{2} - \binom{15}{2} - \binom{12}{2} - \binom{10}{2} = 666 - 105 - 66 - 45 = 450$.
Երկու մեթոդներն էլ տալիս են նույն պատասխանը։ \checkmark

% ══════════════════════════════════════════════════════════════════
\newpage
\section{Խնդիր 5 — 700-ից փոքր բնական թվեր՝ բաժանվող 5-ի և/կամ 3-ի}
% ══════════════════════════════════════════════════════════════════

\begin{theorybox}[Բազմապատիկների հաշվումը միջակայքում]
$N$-ից փոքր կամ հավասար $d$-ի դրական բազմապատիկների քանակը
$\left\lfloor \frac{N}{d} \right\rfloor$ է (ամբողջ մաս)։

Ավելի ճշգրիտ, $\{1, 2, \ldots, N\}$ բազմության մեջ $d$-ի բազմապատիկներն են
$d, 2d, 3d, \ldots, \lfloor N/d \rfloor \cdot d$,
ուստի դրանց քանակը ճշգրիտ $\lfloor N/d \rfloor$ է։

\medskip
\textbf{«$N$-ից փոքր» ընդդեմ «$N$-ից ոչ մեծ».}\;
$N$-ից փոքր բնական թվերն են $\{1, 2, \ldots, N-1\}$։
Այս բազմության մեջ $d$-ի բազմապատիկների քանակը
$\left\lfloor \frac{N-1}{d} \right\rfloor$ է։
\end{theorybox}

\begin{problembox}
700-ից փոքր քանի՞ բնական թիվ կա, որ բաժանվում է 5-ի։
Դրանցից քանիսն են բաժանվում 3-ի։
Քանիսն են բաժանվում և՛ 3-ի, և՛ 5-ի։
\end{problembox}

\subsection*{Լուծում}

Դիտարկենք $S = \{1, 2, 3, \ldots, 699\}$ բազմությունը (700-ից փոքր բնական թվեր)։

\paragraph{(a) 5-ի բաժանվողները.}
$S$-ում 5-ի բազմապատիկներն են $5, 10, 15, \ldots, 695$։
Նման բազմապատիկների քանակն է․
\[
  \left\lfloor \frac{699}{5} \right\rfloor = \lfloor 139.8 \rfloor = 139.
\]

\paragraph{(b) 3-ի բաժանվողները (700-ից փոքր բոլոր բնական թվերի մեջ).}
(Այստեղ «դրանցից»-ը վերաբերում է 700-ից փոքր բոլոր բնական թվերին, ոչ միայն (a) մասի 5-ի բազմապատիկներին։)
$S$-ում 3-ի բազմապատիկներն են $3, 6, 9, \ldots, 699$։
Նման բազմապատիկների քանակն է․
\[
  \left\lfloor \frac{699}{3} \right\rfloor = 233.
\]

\paragraph{(c) Եվ՛ 3-ի, և՛ 5-ի բաժանվողները.}
Թիվը, որը բաժանվում է և՛ 3-ի, և՛ 5-ի, բաժանվում է $\text{ԱԸԲ}(3,5) = 15$-ի։
$S$-ում 15-ի բազմապատիկներն են $15, 30, 45, \ldots, 690$։
Նման բազմապատիկների քանակն է․
\[
  \left\lfloor \frac{699}{15} \right\rfloor = \lfloor 46.6 \rfloor = 46.
\]

\begin{answerbox}
\centering
5-ի բաժանվող՝ $\boxed{139}$։ \qquad
3-ի բաժանվող՝ $\boxed{233}$։ \qquad
Եվ՛ 3-ի, և՛ 5-ի բաժանվող՝ $\boxed{46}$։
\end{answerbox}

% ══════════════════════════════════════════════════════════════════
\newpage
\section{Խնդիր 10 — 7-նիշ թվեր՝ զույգության փոփոխմամբ}
% ══════════════════════════════════════════════════════════════════

\begin{theorybox}[Արտադրյալի կանոնը թվանշանների հաշվման համար]
Երբ հաշվում ենք թվանշանների սահմանափակումներով թվեր, արտադրյալի կանոնը
կիրառվում է դիրք առ դիրք։

\medskip
\textbf{Կարևոր նկատառում.}\;
$n$-նիշ թվի \emph{առաջին թվանշանը} (ձախից առաջինը) \textbf{չի կարող լինել 0},
քանի որ դա կնվազեցնի թվի կարգը։ Այս սահմանափակումը պետք է միշտ հաշվի առնել առանձին։

\medskip
\textbf{Զույգ թվանշաններ.} $\{0, 2, 4, 6, 8\}$ — հինգ հատ։\\
\textbf{Կենտ թվանշաններ.} $\{1, 3, 5, 7, 9\}$ — հինգ հատ։
\end{theorybox}

\begin{problembox}
Գտեք այն 7-նիշ թվերի քանակը, որտեղ զույգ համարի դիրքերում կենտ թվանշաններ են,
իսկ կենտ համարի դիրքերում՝ զույգ թվանշաններ։
\end{problembox}

\subsection*{Լուծում}

Դիրքերը համարակալենք $1, 2, 3, 4, 5, 6, 7$ ձախից աջ։

\begin{center}
\renewcommand{\arraystretch}{1.4}
\begin{tabular}{ccccccc}
\toprule
Դիրք 1 & Դիրք 2 & Դիրք 3 & Դիրք 4 & Դիրք 5 & Դիրք 6 & Դիրք 7 \\
\midrule
\textit{կենտ} & \textit{զույգ} & \textit{կենտ} & \textit{զույգ} & \textit{կենտ} & \textit{զույգ} & \textit{կենտ} \\
(համար) & (համար) & (համար) & (համար) & (համար) & (համար) & (համար) \\
\midrule
զույգ թիվ & կենտ թիվ & զույգ թիվ & կենտ թիվ & զույգ թիվ & կենտ թիվ & զույգ թիվ \\
\bottomrule
\end{tabular}
\end{center}

\paragraph{Զույգ դիրքեր (2, 4, 6).}
Պետք է պարունակեն \textbf{կենտ} թվանշաններ $\{1, 3, 5, 7, 9\}$ բազմությունից։
Յուրաքանչյուր դիրք ունի $5$ ընտրություն, ընդհանուր՝ $5^3$ համակցություն։

\paragraph{Կենտ դիրքեր (1, 3, 5, 7).}
Պետք է պարունակեն \textbf{զույգ} թվանշաններ $\{0, 2, 4, 6, 8\}$ բազմությունից։
Յուրաքանչյուր դիրք ունի $5$ ընտրություն \emph{ընդհանուր առմամբ}, \textbf{բացառությամբ}
դիրք 1-ի (առաջին թվանշան), որը \textbf{չի կարող լինել 0}։

\begin{itemize}[nosep]
  \item Դիրք 1: $\{2, 4, 6, 8\}$ — \textbf{4 ընտրություն}։
  \item Դիրքեր 3, 5, 7: $\{0, 2, 4, 6, 8\}$ — $5$ ընտրություն յուրաքանչյուրը, ընդհանուր՝ $5^3$։
\end{itemize}

\paragraph{Արտադրյալի կանոնի կիրառում.}
Քանի որ ընտրությունները յուրաքանչյուր դիրքում անկախ են․
\[
  \underbrace{4}_{\text{դիրք 1}}
  \;\times\;
  \underbrace{5^3}_{\text{դիրքեր 2, 4, 6}}
  \;\times\;
  \underbrace{5^3}_{\text{դիրքեր 3, 5, 7}}
  \;=\; 4 \times 125 \times 125
  \;=\; 4 \times 15{,}625
  \;=\; 62{,}500.
\]

\textbf{Օրինակ.}\; $2\,143\,658$-ը վավեր 7-նիշ թիվ է (առաջին թվանշանը $\neq 0$)։ Բայց $0\,143\,658$-ը վավեր չէ, քանի որ այն կլիներ 6-նիշ թիվ։ Առաջին թվանշանի սահմանափակումը հեռացնում է բոլոր հնարավոր վավեր տողերի ճիշտ $\frac{1}{5}$-ը (քանի որ 1-ին դիրքում զույգ թվանշաններից $\{0,2,4,6,8\}$ մեկը՝ 0-ն, արգելված է)։

\begin{answerbox}
\centering
Այդպիսի 7-նիշ թվերի քանակն է $\boxed{62{,}500}$։
\end{answerbox}

% ══════════════════════════════════════════════════════════════════
\newpage
\part*{Բաժին 2. Տեղափոխություններ և Դասավորություններ}
\addcontentsline{toc}{part}{Բաժին 2. Տեղափոխություններ և Դասավորություններ}
% ══════════════════════════════════════════════════════════════════

% ══════════════════════════════════════════════════════════════════
\section{Խնդիր 5 — Տղաներ և աղջիկներ շարքում}
% ══════════════════════════════════════════════════════════════════

\begin{theorybox}[Դիրքային սահմանափակումներով տեղափոխություններ]
Երբ շարքում որոշ դիրքեր սահմանափակված են (օրինակ՝ «առաջին և վերջին
դիրքերում պետք է լինեն որոշակի տիպի մարդիկ»), մենք օգտագործում ենք հետևյալ ռազմավարությունը․

\begin{enumerate}[nosep]
  \item \textbf{Նախ լրացրեք սահմանափակված դիրքերը}՝
        հաշվելով դրանց ընտրության հնարավորությունները։
  \item \textbf{Լրացրեք մնացած դիրքերը} մնացած մարդկանցով,
        ովքեր կարող են կանգնել ցանկացած կարգով։
  \item \textbf{Բազմապատկեք} քանակները (արտադրյալի կանոն)։
\end{enumerate}

\medskip
Հիշեցնենք, որ $n$ տարբեր օբյեկտները շարքում դասավորելու եղանակների քանակը
$n! = n \cdot (n-1) \cdot (n-2) \cdots 2 \cdot 1$ է։
\end{theorybox}

\begin{problembox}
Քանի՞ եղանակով կարող են 6 տղա և 7 աղջիկ կանգնել շարքում այնպես,
որ առաջին և վերջին դիրքերում լինեն տղաներ։
\end{problembox}

\subsection*{Լուծում}

Ընդհանուր կա $6 + 7 = 13$ մարդ։

\begin{center}
$\boxed{T}\;\_\;\_\;\_\;\cdots\;\_\;\_\;\boxed{T}$
\end{center}

\paragraph{Քայլ 1. Լրացնել սահմանափակված դիրքերը.}
\begin{itemize}[nosep]
  \item \textbf{Առաջին դիրք.} պետք է լինի տղա։
        Ընտրելու համար կա 6 տղա, ուստի $6$ հնարավորություն։
  \item \textbf{Վերջին դիրք.} պետք է լինի տղա (առաջինից տարբեր)։
        Մնացել է $6 - 1 = 5$ տղա, ուստի $5$ հնարավորություն։
\end{itemize}

\paragraph{Քայլ 2. Լրացնել մնացած դիրքերը.}
Երկու ծայրերում տղաներ տեղադրելուց հետո մնում է $13 - 2 = 11$ մարդ
(4 տղա և 7 աղջիկ)։ Այս 11 մարդիկ կարող են կանգնել մնացած
11 դիրքերում (դիրքեր 2-ից 12) ցանկացած կարգով․
\[
  11! \text{ եղանակով։}
\]

\paragraph{Քայլ 3. Կիրառել արտադրյալի կանոնը.}
Յուրաքանչյուր քայլի ընտրությունները անկախ են, ուստի․
\[
  6 \times 5 \times 11!
  \;=\; 30 \times 11!.
\]

\begin{answerbox}
\centering
Դասավորությունների ընդհանուր քանակն է $\boxed{30 \times 11! = 30 \times 39{,}916{,}800 = 1{,}197{,}504{,}000}$։
\end{answerbox}

% ══════════════════════════════════════════════════════════════════
\newpage
\section{Խնդիր 10 — Ռամսեյի թեորեմ \texorpdfstring{$R(3,3) = 6$}{R(3,3) = 6}}
% ══════════════════════════════════════════════════════════════════

\begin{theorybox}[Դիրիխլեի սկզբունք և Ռամսեյի տեսություն]
\textbf{Դիրիխլեի սկզբունք (Աղավնիների բների սկզբունք).}\;
Եթե $n$ առարկաներ տեղադրվում են $k$ արկղերում և $n > k$, ապա առնվազն
մեկ արկղ պարունակում է մեկից ավելի առարկա։

Ավելի ընդհանուր. եթե $n$ առարկաներ տեղադրվում են $k$ արկղերում, ապա առնվազն
մեկ արկղ պարունակում է $\lceil n/k \rceil$ առարկա։

\medskip
\textbf{Ռամսեյի տեսություն} (ոչ ֆորմալ)․
Ցանկացած բավականաչափ մեծ կառուցվածքում պետք է հայտնվի կարգավորվածություն։
Ռամսեյի թիվը՝ $R(s,t)$-ն, այն ամենափոքր $n$-ն է, որի դեպքում $K_n$ լրիվ գրաֆի
կողերի ցանկացած 2-ներկում (ասենք, կարմիր և կապույտ) պարունակում է
կամ կարմիր $K_s$, կամ կապույտ $K_t$։

\medskip
Դասական արդյունքը $R(3,3) = 6$-ն է․ ցանկացած 6 հոգանոց խմբում միշտ կան
3 փոխադարձ ընկերներ կամ 3 փոխադարձ անծանոթներ։
\end{theorybox}

\begin{problembox}
Ապացուցեք, որ ցանկացած 6 հոգանոց խմբում կան կամ 3 փոխադարձ ընկերներ,
կամ 3 փոխադարձ անծանոթներ։
\end{problembox}

\subsection*{Լուծում (Ապացույց)}

Մենք մոդելավորում ենք իրավիճակը որպես $K_6$ լրիվ գրաֆի 2-ներկում։
Յուրաքանչյուր մարդ գագաթ է, և յուրաքանչյուր կող ներկված է \textbf{կարմիր}
(ընկերներ) կամ \textbf{կապույտ} (անծանոթներ)։
Պետք է ցույց տանք, որ գոյություն ունի միագույն եռանկյուն։

\paragraph{Քայլ 1. Կիրառել Դիրիխլեի սկզբունքը.}
Ընտրեք որևէ $A$ գագաթ։ $A$ գագաթը միացված է մյուս 5 գագաթներին
5 կողերով, որոնցից յուրաքանչյուրը կարմիր է կամ կապույտ։
Ըստ Դիրիխլեի սկզբունքի (5 կող 2 գույնի համար),
$A$-ից դուրս եկող կողերից առնվազն $\lceil 5/2 \rceil = 3$-ը ունեն նույն գույնը։

Առանց ընդհանրությունը խախտելու (WLOG), ենթադրենք $A$-ն ունի առնվազն 3 \textbf{կարմիր}
կող։ Դիցուք այդ կարմիր կողերի մյուս ծայրակետերն են $B$, $C$, $D$-ն։
(3 կապույտ կողերի դեպքը համաչափ է՝ փոխելով «ընկեր» և «անծանոթ» տերմինները։)

\begin{center}
\begin{tikzpicture}[>=Stealth, font=\small,
  vnode/.style={circle, draw, thick, minimum size=7mm, inner sep=0pt}]
  % Vertices in a circle
  \node[vnode, fill=setA!20] (A) at (0:0) {$A$};
  \node[vnode, fill=setB!15] (B) at (60:2)  {$B$};
  \node[vnode, fill=setB!15] (C) at (120:2) {$C$};
  \node[vnode, fill=setB!15] (D) at (180:2) {$D$};
  \node[vnode] (E) at (240:2) {$E$};
  \node[vnode] (F) at (300:2) {$F$};

  % Red edges from A
  \draw[very thick, setB] (A) -- (B);
  \draw[very thick, setB] (A) -- (C);
  \draw[very thick, setB] (A) -- (D);
  % Blue edges from A
  \draw[very thick, setA, dashed] (A) -- (E);
  \draw[very thick, setA, dashed] (A) -- (F);

  % Legend
  \node[font=\scriptsize] at (3.2, 1.0) {\textcolor{setB}{\rule{8pt}{2pt}} կարմիր (ընկերներ)};
  \node[font=\scriptsize] at (3.2, 0.5) {\textcolor{setA}{- - -} կապույտ (անծանոթներ)};

  % Brace for the 3 red neighbors
  \draw[decorate, decoration={brace, amplitude=5pt}, thick, setB!70!black]
    (60:2.7) arc (60:180:2.7);
  \node[font=\scriptsize, setB!70!black] at (120:3.3) {$\ge 3$ նույն գույնի};
\end{tikzpicture}
\end{center}

\paragraph{Քայլ 2. Զննել $B$, $C$, $D$-ի միջև եղած կողերը.}
Դիտարկենք $BC$, $BD$, $CD$ երեք կողերը։
Կա երկու դեպք․

\textbf{Դեպք 1.} $BC$, $BD$, $CD$ կողերից առնվազն մեկը \textbf{կարմիր} է։\\
Ենթադրենք (WLOG) $BC$-ն է կարմիր։ Այդ դեպքում $A$, $B$, $C$-ն կազմում են կարմիր եռանկյուն․
$AB$, $AC$ և $BC$ կողերը բոլորը կարմիր են։
Այսպիսով, մենք ունենք 3 փոխադարձ \textbf{ընկերներ}։

\textbf{Դեպք 2.} $BC$, $BD$, $CD$ կողերից ոչ մեկը կարմիր չէ, այսինքն՝ բոլոր երեքը
\textbf{կապույտ} են։\\
Այդ դեպքում $B$, $C$, $D$-ն կազմում են կապույտ եռանկյուն․ $BC$, $BD$, $CD$ կողերը
բոլորը կապույտ են։
Այսպիսով, մենք ունենք 3 փոխադարձ \textbf{անծանոթներ}։

\medskip
Երկու դեպքում էլ գոյություն ունի միագույն եռանկյուն։ \qed

\paragraph{Ինչու 5 հոգին բավարար չէ.}
Ցույց տալու համար, որ $R(3,3) > 5$, մենք ներկայացնում ենք $K_5$-ի 2-ներկում
առանց միագույն եռանկյունու․ դասավորեք 5 գագաթները կանոնավոր հնգանկյան տեսքով
և ներկեք հնգանկյան կողերը կարմիր, իսկ անկյունագծերը՝ կապույտ (կամ հակառակը)։
Կարելի է ստուգել, որ ոչ մի եռանկյուն միագույն չէ։

\begin{center}
\begin{tikzpicture}[font=\small,
  vnode/.style={circle, draw, thick, minimum size=7mm, inner sep=0pt, fill=white}]
  % 5 vertices in a pentagon
  \foreach \i/\lab in {1/1, 2/2, 3/3, 4/4, 5/5} {
    \node[vnode] (v\i) at ({90+(\i-1)*72}:1.5) {$\lab$};
  }
  % Pentagon edges (red)
  \foreach \i [evaluate=\i as \j using {int(mod(\i,5)+1)}] in {1,...,5} {
    \draw[very thick, setB] (v\i) -- (v\j);
  }
  % Diagonals (blue, dashed)
  \draw[very thick, setA, dashed] (v1) -- (v3);
  \draw[very thick, setA, dashed] (v1) -- (v4);
  \draw[very thick, setA, dashed] (v2) -- (v4);
  \draw[very thick, setA, dashed] (v2) -- (v5);
  \draw[very thick, setA, dashed] (v3) -- (v5);

  \node[font=\footnotesize] at (0, -2.2) {$K_5$: կողերը \textcolor{setB}{կարմիր}, անկյունագծերը \textcolor{setA}{կապույտ} — չկա միագույն եռանկյուն};
\end{tikzpicture}
\end{center}

\begin{answerbox}
\centering
Ցանկացած 6 հոգանոց խմբում միշտ կան 3 փոխադարձ ընկերներ
կամ 3 փոխադարձ անծանոթներ։ Սա ապացուցում է, որ $R(3,3) = 6$։
\end{answerbox}

% ══════════════════════════════════════════════════════════════════
\newpage
\section{Խնդիր 12 — Բաժանելիություն \texorpdfstring{$7, 77, 777, \ldots$}{7, 77, 777, ...} հաջորդականությունում}
% ══════════════════════════════════════════════════════════════════

\begin{theorybox}[Դիրիխլեի սկզբունքը բաժանելիության համար]
Դիրիխլեի սկզբունքի հզոր կիրառություն թվերի տեսությունում․

\medskip
Եթե մենք ունենք $n+1$ ամբողջ թվեր և դիտարկում ենք դրանց մնացորդները $n$-ի վրա բաժանելիս,
ապա ըստ Դիրիխլեի սկզբունքի, դրանցից առնվազն երկուսը ունեն
նույն մնացորդը։ Հետևաբար, նրանց \emph{տարբերությունը} բաժանվում է $n$-ի։

\medskip
\textbf{Repunit-անման հաջորդականություններ.}\;
$7, 77, 777, 7777, \ldots$ հաջորդականությունը կարող է գրվել փակ տեսքով։
$k$-րդ անդամը՝
\[
  a_k = \underbrace{77\ldots7}_{k \text{ հատ յոթ}}
      = 7 \cdot \underbrace{11\ldots1}_{k}
      = 7 \cdot \frac{10^k - 1}{9}.
\]
\end{theorybox}

\begin{problembox}
Ապացուցեք, որ $7,\; 77,\; 777,\; 7777,\; \ldots$ հաջորդականության
առաջին 2014 անդամների մեջ կա մեկը, որը բաժանվում է 2013-ի։
\end{problembox}

\subsection*{Լուծում (Ապացույց)}

Դիցուք $a_k = \underbrace{77\ldots7}_{k}$ որտեղ $k = 1, 2, \ldots, 2014$։

Դիտարկենք 2014 հատ մնացորդները՝ $r_k = a_k \bmod 2013$ երբ $k = 1, 2, \ldots, 2014$։

\paragraph{Դեպք 1. Որևէ $r_k = 0$.}
Եթե $r_k = 0$ որևէ $k$-ի համար, ապա $a_k$-ն բաժանվում է 2013-ի, և խնդիրը լուծված է։

\paragraph{Դեպք 2. Ոչ մի $r_k = 0$.}
Այդ դեպքում յուրաքանչյուր $r_k \in \{1, 2, \ldots, 2012\}$։
Մենք ունենք 2014 մնացորդներ, որոնք ընդունում են արժեքներ 2012 տարր ունեցող բազմությունից։
Ըստ \textbf{Դիրիխլեի սկզբունքի}, գոյություն ունեն $i < j$ ինդեքսներ
($1 \le i < j \le 2014$) այնպես, որ $r_i = r_j$, այսինքն՝
\[
  a_j \equiv a_i \pmod{2013}.
\]
Հետևաբար $2013 \mid (a_j - a_i)$։

\paragraph{$a_j - a_i$-ի վերլուծություն.}
Հաշվենք տարբերությունը․
\[
  a_j - a_i
  = \underbrace{77\ldots7}_{j} - \underbrace{77\ldots7}_{i}
  = \underbrace{77\ldots7}_{j-i} \cdot 10^i
  = a_{j-i} \cdot 10^i.
\]
Ուստի $2013 \mid a_{j-i} \cdot 10^i$։

\textbf{Օրինակ.}\; $\underbrace{77\ldots7}_{5} - \underbrace{77\ldots7}_{3} = 77777 - 777 = 77000 = 77 \times 1000 = \underbrace{77\ldots7}_{2} \times 10^3$։

Այժմ, $2013 = 3 \times 11 \times 61$։ Քանի որ $\gcd(10, 2013) = 1$
(քանի որ 2013-ը չունի 2 կամ 5 բաժանարարներ), հետևում է, որ $\gcd(10^i, 2013) = 1$։

Հետևաբար $2013 \mid a_{j-i}$ (ըստ Գաուսի լեմմայի կամ պարզապես փոխադարձ պարզության)։

Քանի որ $1 \le j - i \le 2013 < 2014$, ապա $a_{j-i}$ անդամը իրապես գտնվում է
հաջորդականության առաջին 2014 անդամների մեջ։

\medskip
Երկու դեպքում էլ, գոյություն ունի անդամ $\{a_1, a_2, \ldots, a_{2014}\}$ բազմության մեջ,
որը բաժանվում է 2013-ի։ \qed

\begin{answerbox}
\centering
Ըստ Դիրիխլեի սկզբունքի, $a_1, a_2, \ldots, a_{2014}$ թվերի մեջ
կա առնվազն մեկը, որը բաժանվում է $2013$-ի։
\end{answerbox}

\end{document}